% Problem dossier: VI.02.P8
\subsection*{Problem VI.02.P8 --- Hyperplane arrangements and a dense complement}
\label{prob:vi02-p8}

\paragraph{Problem.}
Let \(\ell_1,\dots,\ell_r\in k[x_1,\dots,x_n]\) be nonzero affine-linear polynomials over
an infinite field \(k\), and let
\[
H_i=V(\ell_i).
\]
Prove
\[
\mathbb A_k^n\setminus\bigcup_{i=1}^rH_i
=D(\ell_1\cdots\ell_r),
\]
and deduce that this complement is open and dense.

\ifdefined\FullProblemDossiers
\paragraph{Why this problem matters.}
Finite arrangements of forbidden algebraic conditions occur constantly in geometry.  The
Zariski viewpoint compresses many exclusions into one nonvanishing polynomial.

\paragraph{Guided hints.}
Use \(V(fg)=V(f)\cup V(g)\), then take complements.  The product of nonzero linear
polynomials is a nonzero polynomial because the polynomial ring is a domain.

\paragraph{Solution.}
By repeated use of the zero-locus union identity,
\[
V(\ell_1\cdots\ell_r)=\bigcup_{i=1}^rV(\ell_i)=\bigcup_{i=1}^rH_i.
\]
Taking complements gives the displayed identity.  The complement is therefore principal
open.  Since each \(\ell_i\neq0\) and the polynomial ring is a domain, their product is
nonzero.  Over an infinite field every nonzero principal open is dense.

\paragraph{What happened mathematically.}
Finitely many equations to avoid combine multiplicatively into one equation whose
nonvanishing describes the generic region.

\paragraph{Extensions.}
Apply this to the coordinate hyperplanes and recover the dense torus \((k^\times)^n\).

\paragraph{Provenance.}
New canonical problem designed for VI/02.
\fi
